\documentclass[11pt,a4paper]{article}

\usepackage[T1]{fontenc}
\usepackage[utf8]{inputenc}
\usepackage{lmodern}
\usepackage{microtype}
\usepackage[a4paper,margin=1in]{geometry}
\usepackage{amsmath}
\usepackage{amssymb}
\usepackage{amsthm}
\usepackage{mathtools}
\usepackage{booktabs}
\usepackage{pgfplots}
\usepackage{xcolor}
\usepackage[colorlinks=true,linkcolor=blue!55!black,citecolor=blue!55!black,
            urlcolor=blue!55!black,breaklinks=true]{hyperref}

\newtheorem{theorem}{Theorem}
\newtheorem{lemma}{Lemma}
\newtheorem{corollary}{Corollary}
\theoremstyle{remark}
\newtheorem{remark}{Remark}

\newcommand{\id}{\mathrm{id}}
\newcommand{\Tr}{\operatorname{Tr}}
\newcommand{\Ad}{\operatorname{Ad}}
\newcommand{\cD}{\mathcal{D}}
\newcommand{\cE}{\mathcal{E}}
\newcommand{\cN}{\mathcal{N}}
\newcommand{\cT}{\mathcal{T}}
\newcommand{\cU}{\mathcal{U}}
\newcommand{\cL}{\mathcal{L}}
\newcommand{\diamondnorm}[1]{\left\lVert #1\right\rVert_{\diamond}}
\newcommand{\tracenorm}[1]{\left\lVert #1\right\rVert_{1}}
\newcommand{\SWAP}{\mathrm{SWAP}}
\pgfplotsset{compat=1.16}

\hypersetup{
  pdftitle={Exact directional signalling for Ising, partial-SWAP, and XY unitary families},
  pdfauthor={Richard Astbury},
  pdfkeywords={quantum signalling; diamond norm; partial SWAP; Ising interaction;
               quantum channels; semidefinite programming}
}

\title{\bfseries Exact directional signalling for Ising,
partial-SWAP, and XY unitary families}

\author{Richard Astbury\\
\normalsize Independent Researcher\thanks{Contact:
\href{mailto:richard.astbury@gmail.com}{richard.astbury@gmail.com}.
Developed with the assistance of AI tools. The author is solely responsible
for all claims, proofs, code, and errors.}}

\date{Draft of 30 July 2026}

\begin{document}
\maketitle

\begin{abstract}
For a bipartite unitary $U$ on $A\otimes B$, directional signalling from
$B$ to $A$ is the diamond-norm distance
\[
 \delta_A(U)=\inf_{\cE\ {\rm CPTP}}
 \diamondnorm{\Tr_B\!\circ\Ad_U-\cE\circ\Tr_B}.
\]
This quantity was introduced by Barsse, Perinotti, Tosini, and Vaglini.
Exact values were subsequently obtained for the CNOT and SWAP gates, but
to our knowledge continuous interacting families have not been evaluated.
For the two-qubit Ising family
$U_\theta=\exp(-i\theta Z\otimes Z)$, we first prove
$\delta_A(U_\theta)=|\sin 2\theta|$ and identify an optimal dephasing channel.
This elementary baseline fixes conventions and extends the known CNOT
endpoint.  Our first main result concerns the partial-SWAP family on
$\mathbb{C}^d\otimes\mathbb{C}^d$,
$U_\phi=\cos\phi\,I-i\sin\phi\,\SWAP$, where unitary covariance reduces the
optimization to one depolarizing parameter $\lambda$.  We evaluate the
fixed-$\lambda$ diamond norm in closed form and thereby reduce the full
problem to a scalar convex minimization.  In the weak-interaction regime
$0\leq\sin\phi\leq(d^2-3)/(d^2-1)$, the result simplifies to
$\delta_A(U_\phi)=2\sin\phi$; beyond this dimension-dependent threshold the
optimal effective channel moves away from the identity, reaching the known
SWAP value $2(1-1/d^2)$.
Our second main result treats the less symmetric XY line
$e^{-i\theta(XX+YY)}$.  Axial covariance gives an exact
two-parameter min--max formula.  Writing
$S=\sin2\theta$, the weak branch is $\sin4\theta$, while near iSWAP the value is
$S+S^2/2$, reaching the global two-qubit ceiling $3/2$ at iSWAP.  In the intervening branch the value is
the unique physical root of an explicit quartic.
The proofs include matching operational witnesses: explicit input states on
the closed branches, and a convex--concave saddle on the intermediate one.
We also distinguish
these diamond-norm results from Hilbert--Schmidt operator-entanglement
measures and provide exact-arithmetic computational certificates.
\end{abstract}

\noindent\textbf{Keywords:} directional signalling; diamond norm;
partial SWAP; iSWAP; Ising interaction; covariant quantum channels.

\section{Introduction}

An interacting microscopic evolution need not induce autonomous dynamics on
a subsystem.  Given a unitary channel $\Ad_U$ on $A\otimes B$, one may ask how
well the reduced output on $A$ can be simulated from the reduced input on $A$
alone.  The operational answer is
\begin{equation}
 \delta_A(U)
 =\inf_{\cE\in\mathrm{CPTP}(A)}
   \diamondnorm{\Tr_B\circ\Ad_U-\cE\circ\Tr_B}.
 \label{eq:defect}
\end{equation}
The optimization includes arbitrary $AB$ correlations and arbitrary reference
systems through the diamond norm.

Equation~\eqref{eq:defect} is not a new measure.  After exchanging subsystem
labels, it is exactly the directional-signalling quantity introduced in
Ref.~\cite{barsse2024}.  In their notation, signalling from $A$ to $B'$ is
\[
 S_{A\to B'}(\cU)
 =\inf_{\mathcal C\ {\rm CPTP}}
 \diamondnorm{\Tr_{A'}\circ\cU-\mathcal C\circ\Tr_A}.
\]
Thus Eq.~\eqref{eq:defect} is $S_{B\to A'}(\Ad_U)$: exchange their
$A,A'$ with our hidden $B,B'$ and their $B,B'$ with our visible $A,A'$.
Both definitions use the standard, unhalved diamond norm, so there is no
normalization factor.  Ref.~\cite{barsse2024} gives
$S(\mathrm{CNOT})\leq1$ as part of proving a strict separation from causal
influence.  The follow-up Ref.~\cite{barsse2025} proves the matching lower
bound and hence the exact equality $S(\mathrm{CNOT})=1$; it also treats SWAP,
including parallel and asymptotic uses, and gives the qudit SWAP value used
below.  It further proves the universal bound
\begin{equation}
 \delta_A(U)\leq2\left(1-\frac1{d^2}\right)
 \label{eq:universal-swap-bound}
\end{equation}
for equal local dimension $d$, with equality at SWAP, and uses covariance to
identify the completely depolarizing effective channel there.  Our
partial-SWAP endpoint uses that precedent; the contribution is the continuous
family and its threshold.  Earlier work characterized causal and semicausal
operations structurally~\cite{beckman2001,eggeling2002}.  In the present
notation, $\delta_A(U)=0$ is exactly semicausality from $B$ to $A$: the
reduced output on $A$ then factors through the reduced input on $A$.
Continuity of Stinespring dilations supplies general non-sharp
bounds~\cite{ksw2008}.

The present contribution is deliberately narrower than a new causal measure.
The Ising calculation is a transparent baseline that fixes the subsystem and
diamond-norm conventions while extending the known CNOT endpoint to its
continuous Cartan axis~\cite{zhang2003}.  The two substantive results are the
arbitrary-dimensional partial-SWAP reduction and the complete XY/iSWAP phase
diagram.  Partial SWAP is used in homogenization and collision models
\cite{ziman2002,scarani2002,ciccarello2022} and here exhibits the
dimension-dependent
transition $\sin\phi=(d^2-3)/(d^2-1)$.  The less symmetric XY line has three
analytic regimes, an algebraic quartic in the middle, and an iSWAP endpoint
that saturates the global two-qubit bound.  A systematic search of the
directional-signalling, semicausal-channel, collision-model, and two-qubit-gate
literatures found no previous evaluation of Eq.~\eqref{eq:defect} for these
continuous families.  The search scope and attribution boundary are recorded
in Appendix~\ref{app:literature}; this remains a literature assessment rather
than a theorem.

The unifying phenomenon is not merely the availability of three formulas.
As the coupling grows, the best reduced description can move from the
identity to a deliberately noisy channel, and the input witnessing the
unavoidable error can change at the same time.  The partial-SWAP and XY
thresholds make this active-set transition exact.  In collision-model
language, they identify when retaining the incoming visible state ceases to
be the best one-step autonomous approximation.

The distinction from operator entanglement is important.  Hilbert--Schmidt
distances and the linear entropy of a vectorized unitary are useful algebraic
diagnostics~\cite{zanardi2001}, but they do not directly give operational
channel distinguishability.  Every result below is stated in diamond norm and
is accompanied by both an upper bound and a matching witness.
Likewise, tensor-product-structure geometry
\cite{andreadakis2025}, approximate operator-algebra correction
\cite{beny2010}, and exact Lindblad model reduction
\cite{grigoletto2025} address related questions of subsystem structure,
recoverability, and reduced description, but optimize different objects and
do not yield Eq.~\eqref{eq:defect}.  Operational quantum mereology also seeks
dynamically preferred, minimally scrambled subsystems~\cite{zanardi2024}, but
uses short-time scrambling diagnostics rather than worst-case diamond-norm
channel simulation.

\subsection*{Operational picture}

Informally, $A$ is a visible system and $B$ is hidden.  They interact, $B$ is
discarded, and one tries to replace the joint evolution by the best channel
acting on $A$ alone.  The number $\delta_A(U)$ is the unavoidable worst-case
error of that replacement, including inputs correlated with $B$ and with an
external reference.  A zero value means that $A$ admits autonomous reduced
dynamics.  A nonzero value quantifies how strongly information hidden in $B$
can affect the visible output.

Figure~\ref{fig:xy-curve} shows the most intricate example below.  The formulas
change because both the optimal effective channel and the input witnessing its
error change with interaction strength.

\begin{figure}[ht]
\centering
\begin{tikzpicture}
\begin{axis}[
 width=0.82\textwidth,
 height=0.42\textwidth,
 xlabel={interaction angle $\theta$},
 ylabel={unavoidable error $\delta_A$},
 xmin=0, xmax=0.7853982,
 ymin=0, ymax=1.6,
 xtick={0,0.25827052,0.64661551,0.78539816},
 xticklabels={$0$,$\theta_1$,$\theta_2$,$\pi/4$},
 grid=major,
]
\addplot[black!55,densely dotted,thick] coordinates {(0,1.5) (0.7853982,1.5)};
\addplot[blue!65!black,very thick] table[x=theta,y=defect] {xy_curve.dat};
\draw[dashed] (axis cs:0.25827052,0) -- (axis cs:0.25827052,1.6);
\draw[dashed] (axis cs:0.64661551,0) -- (axis cs:0.64661551,1.6);
\node[anchor=west] at (axis cs:0.04,0.22) {weak};
\node[anchor=west] at (axis cs:0.37,1.16) {quartic};
\node[anchor=east] at (axis cs:0.76,1.43) {iSWAP};
\node[anchor=west,black!65] at (axis cs:0.02,1.52) {two-qubit ceiling};
\end{axis}
\end{tikzpicture}
\caption{The exact XY directional-signalling curve.  The middle segment is the
unique physical root of Eq.~\eqref{eq:xy-middle-quartic}.}
\label{fig:xy-curve}
\end{figure}

\section{Preliminaries}

Let $\cL(H)$ denote the operators on a finite-dimensional Hilbert space $H$.
For a unitary $U$, write $\Ad_U(X)=UXU^\dagger$.  We use the unnormalized Choi
operator
\[
 J(\Phi)=\sum_{j,k}\Phi(|j\rangle\langle k|)
                 \otimes |j\rangle\langle k|.
\]
For Hermiticity-preserving $\Phi:\cL(X)\to\cL(Y)$, we use the Watrous
semidefinite-program characterization~\cite{watrous2009,watrous2018}.
In the unnormalized-Choi convention used here, its Hermiticity-preserving
linear primal specialization is
\begin{equation}
 \diamondnorm{\Phi}
 =\max\left\{\Tr[J(\Phi)Y]:
 \begin{array}{l}
  \rho\geq0,\quad \Tr\rho=1,\quad Y=Y^\dagger,\\
  -I_Y\otimes\rho\leq Y\leq I_Y\otimes\rho
 \end{array}\right\}.
 \label{eq:watrous-linear}
\end{equation}
Optimizing first over $Y$ gives the equivalent trace-norm expression
\begin{equation}
 \diamondnorm{\Phi}
 =\max_{\rho\in\mathrm{D}(X)}
 \tracenorm{(I_Y\otimes\sqrt{\rho})J(\Phi)
                     (I_Y\otimes\sqrt{\rho})}.
 \label{eq:watrous}
\end{equation}
The maximizing $\rho$ specifies a physical witness through any purification.
If $\Phi$ is covariant under input $R_g$ and output $S_g$, then
\[
 (\rho,Y)\longmapsto
 \left(\overline{R_g}\rho R_g^{\mathsf T},
 (S_g\otimes\overline{R_g})Y
 (S_g\otimes\overline{R_g})^\dagger\right)
\]
preserves every constraint and the objective in
Eq.~\eqref{eq:watrous-linear}.  Haar averaging this complete feasible pair
therefore produces an invariant optimal density.  This linear averaging, not
averaging $\rho$ inside the nonlinear right-hand side of
Eq.~\eqref{eq:watrous}, is the symmetry argument used below.  The same SDP
also shows that its optimal value is concave in $\rho$: feasible pairs at two
densities can be convexly combined.

We use the maps
\[
 \cT=\Tr_B,\qquad \cN_U=\Tr_B\circ\Ad_U,
\]
so that $\delta_A(U)=\inf_\cE\diamondnorm{\cN_U-\cE\cT}$.
Pre- and post-composition by unitary channels are diamond-norm isometries.

\begin{lemma}[Covariant twirling]
\label{lem:twirl}
Suppose a compact group $G$ has unitary representations $H_g$ on $A\otimes B$
and $G_g$ on $A$ such that
\[
 \cN_U\Ad_{H_g}=\Ad_{G_g}\cN_U,\qquad
 \cT\Ad_{H_g}=\Ad_{G_g}\cT.
\]
Then some optimal effective channel in Eq.~\eqref{eq:defect} is covariant under
$G_g$.
\end{lemma}

\begin{proof}
For a candidate $\cE$, define
$\cE_g=\Ad_{G_g^\dagger}\cE\Ad_{G_g}$.  Covariance gives
\[
 \cN_U-\cE_g\cT
 =\Ad_{G_g^\dagger}(\cN_U-\cE\cT)\Ad_{H_g},
\]
so every $\cE_g$ has the same defect as $\cE$.  The Haar average
$\overline{\cE}=\int_G\cE_g\,dg$ is CPTP and covariant.  Convexity of the
diamond norm implies that its defect is no larger.
\end{proof}

\section{The Ising family}

Let
\[
 U_\theta=e^{-i\theta Z\otimes Z}
          =cI-is\,Z\otimes Z,\qquad c=\cos\theta,\quad s=\sin\theta.
\]

\begin{theorem}[Exact Ising signalling]
\label{thm:ising}
For two qubits,
\begin{equation}
 \delta_A(U_\theta)=|\sin 2\theta|=2|cs|.
 \label{eq:ising-result}
\end{equation}
An optimal effective channel is
\[
 \cE_\theta(X)=c^2X+s^2ZXZ.
\]
\end{theorem}

\begin{proof}
Define
\[
 M(X)=\Tr_B[(I\otimes Z)X],\qquad
 L(X)=-\frac{i}{2}[Z,M(X)].
\]
Expanding $U_\theta XU_\theta^\dagger$ and using cyclicity inside the partial
trace gives
\begin{equation}
 \cN_{U_\theta}-\cE_\theta\cT
 =2cs\,L.
 \label{eq:ising-factor}
\end{equation}
The completely bounded trace norm of $M$ is one: its adjoint is the complete
operator-norm isometry $Y\mapsto Y\otimes Z$.  Also
$K\mapsto-i[Z,K]/2$ has completely bounded trace norm at most one by the
triangle inequality for left and right unitary multiplication.  Hence
$\diamondnorm{L}\leq1$, which proves the upper bound $2|cs|$.

For the lower bound, consider
\[
 \rho_0=|+\rangle\langle+|_A\otimes|0\rangle\langle0|_B,\qquad
 \rho_1=|+\rangle\langle+|_A\otimes|1\rangle\langle1|_B.
\]
They have the same $A$ marginal, so every candidate $\cE$ assigns them the
same simulated output.  Their true reduced outputs differ by $2cs\,Y$, whose
trace norm is $4|cs|$.  The triangle inequality therefore forces at least one
of the two simulation errors to be at least $2|cs|$.  This matches the upper
bound.
\end{proof}

At $\theta=\pi/4$, $U_\theta$ is locally equivalent to CZ and CNOT.
Equation~\eqref{eq:ising-result} then reproduces the exact CNOT value one,
already obtained in Ref.~\cite{barsse2025}.  The continuous
formula, rather than that endpoint, is the result of
Theorem~\ref{thm:ising}.

\section{The partial-SWAP family in dimension \texorpdfstring{$d$}{d}}

Let $A\cong B\cong\mathbb{C}^d$ and consider
\begin{equation}
 U_\phi=cI-is\,\SWAP,\qquad c=\cos\phi,\quad s=\sin\phi,
 \label{eq:partial-swap}
\end{equation}
for $0\leq\phi\leq\pi/2$.  For qubits, the equal-angle Cartan interaction is
this family up to global phase~\cite{zhang2003}:
\[
 e^{-it(XX+YY+ZZ)}=e^{it}e^{-2it\SWAP},\qquad \phi=2t.
\]

\subsection{Reduction to depolarizing channels}

The unitary $U_\phi$ commutes with $V\otimes V$ for every $V\in U(d)$.
Lemma~\ref{lem:twirl} therefore permits an optimal effective channel to be
chosen covariant under all unitary conjugations.  Every such channel is
depolarizing:
\begin{equation}
 \cE_\lambda(X)
 =\lambda X+(1-\lambda)\Tr(X)\frac{I}{d},
 \qquad -\frac{1}{d^2-1}\leq\lambda\leq1.
 \label{eq:depol}
\end{equation}
This interval is precisely the CPTP range.  Direct expansion gives
\[
 \cN_{U_\phi}(X)
 =c^2\Tr_BX+s^2\Tr_AX-ics\,\Tr_B(SX-XS),
\]
where $\Tr_A X$ is relabelled onto the output copy of $A$.

\subsection{The fixed-channel diamond norm}

Define
\begin{align}
 r&=1-\lambda,&
 b&=r^2+d^2\lambda s^2,\nonumber\\
 h_0&=\frac{d(d^2-3)}{d^2-1},&
 h_1&=\frac{2}{d^2-1}.
 \label{eq:general-terms}
\end{align}
For $b>0$, also put
\begin{equation}
 q=r^2(h_0^2-h_1^2)+4b,\quad
 k=1-\frac{r^2h_1^2}{4b},\quad
 z=\frac{rh_0+\sqrt q}{k}.
 \label{eq:general-D}
\end{equation}

\begin{lemma}[Fixed depolarizing channel in dimension $d$]
\label{lem:fixed}
Let
\[
 d_{d,\phi}(\lambda)
 =\diamondnorm{\cN_{U_\phi}-\cE_\lambda\cT}.
\]
Apart from the identity point $b=0$, where the norm vanishes,
\begin{equation}
 d_{d,\phi}(\lambda)=
 \begin{cases}
  z/d, & rh_1z\leq4b,\\[2pt]
  2r(h_0+h_1)/d, & rh_1z>4b.
 \end{cases}
 \label{eq:fixed-norm}
\end{equation}
The quantities $b$ and $k$ are positive and $q$ is nonnegative on the CPTP
interval away from that identity point, so the closed form is well defined.
The lemma is asserted, and proved, on the \emph{whole} admissible interval
$-1/(d^2-1)\leq\lambda\leq1$, including negative $\lambda$; the outer
minimization later uses only $0\leq\lambda\leq1$, so the closed form is
established on a strictly larger range than is subsequently needed.
\end{lemma}

\begin{proof}
The Watrous SDP is invariant under the simultaneous $V\otimes V$ covariance
action.  We average the full linear primal feasible tuple, not the nonlinear
right-hand side of Eq.~\eqref{eq:watrous}.  Affineness preserves feasibility
and objective value, and Schur--Weyl duality then permits an optimal input
density
\[
 \rho=xP_-+yP_+,\qquad r_-x+r_+y=1,
 \quad r_\pm=\frac{d(d\pm1)}2.
\]
Let $u=r_+y$ be its total symmetric-subspace weight.  Put
\[
 \kappa_+=\frac{(d-1)(d+2)}{d+1},\qquad
 \kappa_-=\frac{(d+1)(d-2)}{d-1},\qquad
 \mathcal B=(1+\lambda)^2-4\lambda c^2.
\]
For
\[
 K_\rho=(I\otimes\sqrt\rho)
 J(\cN_{U_\phi}-\cE_\lambda\cT)
 (I\otimes\sqrt\rho),
\]
the block decomposition can be made explicit.  On output--input space
$A_{\rm out}\otimes A_{\rm in}\otimes B_{\rm in}$ define
\[
 |v_a^\pm\rangle=
 \frac{1}{\sqrt{2(d\pm1)}}\sum_{j=1}^d
 |j\rangle\otimes
 (|j,a\rangle\pm|a,j\rangle),\qquad a=1,\ldots,d.
\]
The vectors $v_a^\pm$ span the only subspaces coupled between the symmetric
and antisymmetric input sectors.  Inside
$A_{\rm out}\otimes\operatorname{Sym}^2(\mathbb C^d)$, the
orthogonal complement of $\operatorname{span}\{v_a^+\}$ carries the scalar
$-(1-\lambda)y/d$.  Inside
$A_{\rm out}\otimes\wedge^2(\mathbb C^d)$, the orthogonal
complement of $\operatorname{span}\{v_a^-\}$ carries
$-(1-\lambda)x/d$.  For each $a$, the two-dimensional space
$\operatorname{span}\{v_a^-,v_a^+\}$ carries, in that order, the block
\[
 \begin{pmatrix}
  \dfrac{(d+1)(d-2)}{2d}(1-\lambda)x&\overline\omega\\[5pt]
  \omega&\dfrac{(d-1)(d+2)}{2d}(1-\lambda)y
 \end{pmatrix},
 \qquad
  |\omega|^2=\frac{d^2-1}{4}xy\,\mathcal B .
\]
The phase of $\omega$ depends on the phase convention for $v_a^\pm$ and does
not affect the spectrum.  Using $u=r_+y$, $1-u=r_-x$, its trace and determinant
give the remaining roots below.  Thus the complete spectrum is
\begin{center}
\begin{tabular}{@{}lll@{}}
\toprule
eigenvalue & multiplicity\\
\midrule
$-(1-\lambda)y/d$ & $d(d-1)(d+2)/2$\\
$-(1-\lambda)x/d$ & $d(d+1)(d-2)/2$\\
$\nu_+,\nu_-$ & $d$ each\\
\bottomrule
\end{tabular}
\end{center}
where
\begin{align}
 \nu_++\nu_-&=\frac{1-\lambda}{d^2}
 [\kappa_+u+\kappa_-(1-u)],\nonumber\\
 \nu_+\nu_-&=\frac{u(1-u)}{d^2}
 \left[\frac{(1-\lambda)^2}{d^2}\kappa_+\kappa_- -\mathcal B\right].
 \label{eq:general-roots}
\end{align}
Indeed $\kappa_+\kappa_-=d^2-4$ and
\[
 \mathcal B-\frac{(1-\lambda)^2}{d^2}\kappa_+\kappa_-
 =\frac{4b}{d^2}\geq0,
\]
so the last product is nonpositive even when $\lambda<0$.  Summing all
eigenvalue magnitudes gives
\begin{equation}
 d_{d,\phi}(\lambda)
 =\max_{0\leq u\leq1}\left\{
 \frac{1-\lambda}{d}[\kappa_+u+\kappa_-(1-u)]
 +\sqrt{\frac{(1-\lambda)^2}{d^2}
 [\kappa_+u-\kappa_-(1-u)]^2+4u(1-u)\mathcal B}
 \right\}.
 \label{eq:general-u}
\end{equation}

With $t=2u-1$, elementary algebra turns Eq.~\eqref{eq:general-u} into
\begin{equation}
 \max_{-1\leq t\leq1}\frac1d\left[
 r(h_0+h_1t)
 +\sqrt{r^2(h_0+h_1t)^2+4b(1-t^2)}
 \right].
 \label{eq:general-t}
\end{equation}
Here one uses
$\kappa_+\kappa_-=d^2-4=h_0^2-h_1^2$; this identity accounts for the
apparently different linear and radical coefficients in
Eqs.~\eqref{eq:general-u} and~\eqref{eq:general-t}.
If the maximum is interior, its location is
$t_*=rh_1z/(4b)$ and its value is $z/d$.  Since every factor in $t_*$ is
nonnegative, a maximum at $t=-1$ cannot occur.  The condition $t_*\leq1$
gives the first branch of Eq.~\eqref{eq:fixed-norm}; otherwise the maximum is
at $t=1$ and gives its second branch.  The dimensions of the two scalar
complements are respectively $d(d-1)(d+2)/2$ and
$d(d+1)(d-2)/2$, so the displayed blocks account for all $d^3$ eigenvalues.
The maximizing $\rho$, together with any purification, realizes the matching
lower-bound witness; Eq.~\eqref{eq:watrous-linear} supplies the upper bound.
\end{proof}

\begin{remark}[Status of the general-$d$ spectrum]
\label{rem:generald}
The block decomposition just given is a hand calculation valid for every
$d\geq2$: the vectors $v_a^\pm$ are defined uniformly in $d$, and the stated
multiplicities $d(d-1)(d+2)/2$, $d(d+1)(d-2)/2$, and $d$, $d$ sum to $d^3$,
so no eigenvalue is unaccounted for at any dimension.  The accompanying
script \texttt{tools/verify\_block\_reductions.py} reconstructs the unreduced
physical Choi operator and confirms this spectrum for $d=2,3,4$ only.  Those
checks are therefore a spot check of the algebra above, not the source of the
all-$d$ claim, which rests on the Schur--Weyl argument in the proof.
\end{remark}

\subsection{Optimization and the dimension-dependent threshold}

\begin{theorem}[Partial-SWAP signalling in arbitrary dimension]
\label{thm:partial-swap}
For every $d\geq2$,
\begin{equation}
 \delta_A(U_\phi)
 =\min_{-1/(d^2-1)\leq\lambda\leq1}d_{d,\phi}(\lambda),
 \label{eq:scalar-min}
\end{equation}
with $d_{d,\phi}$ given explicitly by Eq.~\eqref{eq:fixed-norm}.  Moreover,
\begin{equation}
 0\leq\sin\phi\leq\frac{d^2-3}{d^2-1}
 \quad\Longrightarrow\quad
 \lambda_*=1,\qquad
 \delta_A(U_\phi)=2\sin\phi.
 \label{eq:weak-branch}
\end{equation}
For stronger non-SWAP coupling, a minimizer lies in $0<\lambda_*<1$.
Equation~\eqref{eq:scalar-min} is an exact one-dimensional convex
characterization in this regime.  For qubits the interior minimizer is
algebraic: writing $s=\sin\phi$, it satisfies the unsquared stationarity
equation
\begin{equation}
 3(\lambda_*-1+2s^2)
 \left[\sqrt{(1+\lambda_*)^2-4\lambda_*(1-s^2)}
       -\frac{2(1-\lambda_*)}{3}\right]
 =(1+\lambda_*)^2-4\lambda_*(1-s^2).
 \label{eq:qubit-swap-stationarity}
\end{equation}
In particular it is a root of the explicit quartic
\begin{align}
 \mathcal P(\lambda,s)={}&2\lambda^4+(18s^2-8)\lambda^3
 +(45s^4-43s^2+12)\lambda^2\nonumber\\
 &+(36s^6-54s^4+32s^2-8)\lambda
 +5s^4-7s^2+2,\qquad
 \mathcal P(\lambda_*,s)=0.
 \label{eq:qubit-swap-quartic}
\end{align}
Roots introduced by eliminating the square root are discarded using
Eq.~\eqref{eq:qubit-swap-stationarity}.  In particular, positivity of its
right-hand side and of the bracket on the physical branch forces
\begin{equation}
 \lambda_*>1-2s^2.
 \label{eq:qubit-swap-selector}
\end{equation}
The quartic may have an additional root in $(0,1)$, but that root solves the
conjugate equation with the opposite square-root sign.  The endpoint
factorizations
\[
 \mathcal P(0,s)=(s^2-1)(5s^2-2),\qquad
 \mathcal P(1,s)=4s^4(9s^2-1)
\]
for the left-hand side $\mathcal P$ of
Eq.~\eqref{eq:qubit-swap-quartic} make these possible extraneous-root changes
explicit.  For general $d$, Eq.~\eqref{eq:scalar-min} remains the exact
characterization; we do not claim a useful generic radicals expression after
the outer minimization.
At $\phi=\pi/2$,
\begin{equation}
 \lambda_*=0,\qquad
 \delta_A(S)=2\left(1-\frac1{d^2}\right).
 \label{eq:swap-d}
\end{equation}
\end{theorem}

\begin{proof}
Equation~\eqref{eq:scalar-min} follows from covariance and
Lemma~\ref{lem:fixed}.  The function $d_{d,\phi}(\lambda)$ is convex because it
is the diamond norm of a map affine in $\lambda$.  At $\lambda=1$,
\[
 d_{d,\phi}(1)=2s,\qquad
 \left.\frac{d}{d\lambda}d_{d,\phi}(\lambda)\right|_{1^-}
 =s-\frac{d^2-3}{d^2-1}.
\]
Convexity proves Eq.~\eqref{eq:weak-branch}.  The matching invariant density
has equal total symmetric and antisymmetric weights,
\[
 \rho=\frac{dI-S}{d(d^2-1)}.
\]

At $\lambda=0$ and $c>0$, the derivative is
$-2(d^2-1)c^2/d^2<0$.  Convexity excludes negative $\lambda$ from the outer
optimum, while the positive derivative at $\lambda=1$ in the strong regime
places a minimizer in $(0,1)$.  At SWAP the minimizer is $\lambda=0$;
Eq.~\eqref{eq:general-t} is maximized at $t=1/d$, corresponding to
$\rho=I/d^2$, and yields Eq.~\eqref{eq:swap-d}, in agreement with
Ref.~\cite{barsse2025}.

For $d=2$ and $1/3<s<1$, differentiating the interior-maximizer branch of
Eq.~\eqref{eq:fixed-norm} gives
Eq.~\eqref{eq:qubit-swap-stationarity}.  On $0<\lambda<1$ the interior-maximizer branch is automatic, since
\[
 b-\frac{4(1-\lambda)^2}{9}
 =\frac{5(1-\lambda)^2+36\lambda s^2}{9}>0.
\]
Squaring once gives $-4\mathcal P(\lambda,s)=0$.  Equation
\eqref{eq:qubit-swap-selector} records the sign lost in this squaring step.
Conversely, any $\lambda\in(0,1)$ satisfying the unsquared equation has zero
derivative and is therefore globally minimizing by convexity.  There can be
only one such root: two distinct stationary minimizers would make the convex
analytic objective constant on an interval, contrary to its nonconstant
derivative.  Thus Eq.~\eqref{eq:qubit-swap-selector} selects $\lambda_*$ from
any conjugate-sign quartic root.
\end{proof}

\begin{table}[t]
\centering
\caption{Dimension dependence of the partial-SWAP weak threshold and the known
full-SWAP endpoint.}
\label{tab:values}
\begin{tabular}{@{}lll@{}}
\toprule
$d$ & weak threshold for $\sin\phi$ & $\delta_A(S)$\\
\midrule
$2$ & $1/3$ & $3/2$\\
$3$ & $3/4$ & $16/9$\\
$4$ & $13/15$ & $15/8$\\
$d$ & $(d^2-3)/(d^2-1)$ & $2(1-1/d^2)$\\
\bottomrule
\end{tabular}
\end{table}

\begin{figure}[t]
\centering
\begin{tikzpicture}
\begin{axis}[
 width=0.82\textwidth,
 height=0.40\textwidth,
 xlabel={partial-SWAP angle $\phi$},
 ylabel={optimizer parameter},
 xmin=0, xmax=1.5707964,
 ymin=0, ymax=1.04,
 xtick={0,0.3398369,0.7853982,1.5707964},
 xticklabels={$0$,$\arcsin(1/3)$,$\pi/4$,$\pi/2$},
 legend style={at={(0.02,0.04)},anchor=south west},
 grid=major,
]
\addplot[blue!65!black,very thick]
 table[x=phi,y=shrinkage] {partial_swap_phase.dat};
\addlegendentry{channel shrinkage $\lambda_*$}
\addplot[orange!80!black,very thick,dashed]
 table[x=phi,y=symmetric_weight] {partial_swap_phase.dat};
\addlegendentry{symmetric witness weight $u_*$}
\draw[densely dashed] (axis cs:0.3398369,0) -- (axis cs:0.3398369,1.04);
\end{axis}
\end{tikzpicture}
\caption{Optimizer transition for qubit partial SWAP.  The effective channel
remains the identity until $\sin\phi=1/3$ and then becomes progressively more
depolarizing.  Simultaneously, the invariant diamond witness shifts from equal
total symmetric/antisymmetric weight $u_*=1/2$ toward $u_*=3/4$ at SWAP.
The plotted interior values solve the scalar optimization in
Eq.~\eqref{eq:scalar-min}.}
\label{fig:partial-swap-phase}
\end{figure}

\section{The XY/iSWAP line}

Consider
\begin{equation}
 U_\theta=e^{-i\theta(XX+YY)},\qquad 0\leq\theta\leq\frac{\pi}{4},
 \qquad C=\cos2\theta,\quad S=\sin2\theta.
 \label{eq:xy-unitary}
\end{equation}
The zero- and two-excitation states are fixed, while on
$\operatorname{span}\{|01\rangle,|10\rangle\}$ the unitary is
$C I-iS X$.  Equivalently,
\[
 U_\theta=\cos^2\theta\,II-\frac{i}{2}\sin2\theta(XX+YY)
           +\sin^2\theta\,ZZ.
\]

\subsection{Axial covariance and the fixed-channel norm}

Joint Pauli covariance first restricts an optimal effective channel to a Pauli
channel.  The additional axial symmetry generated by $ZI+IZ$ makes its $X$
and $Y$ weights equal.  We therefore write
\begin{equation}
 \cE_{u,w}:\quad (r_x,r_y,r_z)\longmapsto(ur_x,ur_y,wr_z),
 \qquad -1\leq w\leq1,\quad |u|\leq\frac{1+w}{2}.
 \label{eq:xy-channel}
\end{equation}
The corresponding Pauli probabilities are
\[
 q_I=\frac{1+2u+w}{4},\qquad q_X=q_Y=\frac{1-w}{4},\qquad
 q_Z=\frac{1+w-2u}{4}.
\]

\begin{lemma}[XY fixed-channel norm]
\label{lem:xy-fixed}
For fixed $(u,w)$ in Eq.~\eqref{eq:xy-channel}, put
$B=\frac12-A$ and
\[
 \alpha=\frac{A(1-w)}2,\quad \gamma^2=\frac{AB}{2},\quad
 g=C-u,\quad h=S,\quad
 \beta=\frac{B}{2}\left[w-(C^2-S^2)+2iSC\right].
\]
Define
\begin{equation}
 M_3(A)=
 \begin{pmatrix}
 \alpha&\gamma(g+ih)&\gamma(-g+ih)\\
 \gamma(g-ih)&0&\beta\\
 \gamma(-g-ih)&\overline\beta&0
 \end{pmatrix}.
 \label{eq:xy-m3}
\end{equation}
Then
\begin{equation}
 \diamondnorm{\cN_{U_\theta}-\cE_{u,w}\cT}
 =\max_{0\leq A\leq1/2}
 2\left[\alpha+\tracenorm{M_3(A)}\right].
 \label{eq:xy-fixed-norm}
\end{equation}
\end{lemma}

\begin{proof}
Average the complete linear primal diamond-SDP tuple over joint axial
rotations and $X\otimes X$; averaging only the density in
Eq.~\eqref{eq:watrous} would not justify the reduction.  The residual
odd-sector coherence changes sign under the
real-linear involution
\[
 X\longmapsto(Z\otimes I)X^*(Z\otimes I),
 \]
which is a covariance because
$U_\theta(Z\otimes I)=(Z\otimes I)U_\theta^*$.  Averaging this involution leaves
an optimal density of the form
\[
 \rho=A(|00\rangle\langle00|+|11\rangle\langle11|)
 +B(|01\rangle\langle01|+|10\rangle\langle10|).
\]
Inserting $\rho$ in Eq.~\eqref{eq:watrous}, the scaled Choi operator commutes
with the axial charge
\[
 Q_{\rm ax}=\frac12
 (Z_{A_{\rm out}}-Z_{A_{\rm in}}-Z_{B_{\rm in}}),
\]
where the input signs reflect the conjugate covariance action.  Its four
charge sectors have dimensions $1,3,3,1$.  The sector spanned by
$|011\rangle$ contributes the eigenvalue $-\alpha$, while the ordered basis
\[
 |000\rangle,\qquad
 \frac{|101\rangle+|110\rangle}{\sqrt2},\qquad
 \frac{|101\rangle-|110\rangle}{\sqrt2}
\]
gives Eq.~\eqref{eq:xy-m3}.  The symmetry $X^{\otimes3}$ reverses the charge
and carries these sectors to the isolated vector $|100\rangle$ and the
three-dimensional sector formed from $|111\rangle,|010\rangle,|001\rangle$.
The latter is unitarily equivalent after conjugating phases.  This accounts
for all eight dimensions and proves Eq.~\eqref{eq:xy-fixed-norm}.  For
reference, the characteristic polynomial
$x^3+p_2x^2+p_1x+p_0$ of $M_3$ has
\begin{align*}
 p_2&=-\alpha,\\
 p_1&=-\left[2\gamma^2(g^2+h^2)+|\beta|^2\right],\\
 p_0&=\alpha|\beta|^2
 +2\gamma^2(g^2-h^2)\Re\beta
 -4\gamma^2gh\Im\beta.
\end{align*}
\end{proof}

Lemma~\ref{lem:xy-fixed} is an exact finite-dimensional answer along the whole
XY line:
\begin{equation}
 \delta_A(U_\theta)=
 \min_{\substack{-1\leq w\leq1\\|u|\leq(1+w)/2}}
 \max_{0\leq A\leq1/2}
 2\left[\alpha+\tracenorm{M_3(A)}\right].
 \label{eq:xy-minmax}
\end{equation}
Whenever $M_3$ has one nonnegative and two nonpositive eigenvalues,
$\Tr M_3=\alpha$ gives
$2[\alpha+\|M_3\|_1]=4\lambda_{\max}(M_3)$.
Two coupling regimes simplify further.

\subsection{Closed weak and iSWAP branches}

Let $\theta_1=0.258270520262\ldots$ be the unique root in
$(0,\pi/4)$ of
\begin{equation}
 2SC=3C^3+C^2-C-1,
 \label{eq:xy-first-threshold}
\end{equation}
and let $\theta_2=0.646615513406\ldots$ solve
\begin{equation}
 1+C^2=2C(1+S).
 \label{eq:xy-second-threshold}
\end{equation}
Equivalently, $\tan\theta_2$ is the positive root of $t^3+t^2-1=0$.
The middle interval lies inside $(\pi/24,5\pi/24)$.  Indeed, direct
substitution at $\theta=\pi/24$ gives
\[
 2SC-(3C^3+C^2-C-1)
 =1-\frac{11\sqrt2+5\sqrt6}{16}-\frac{\sqrt3}{4}<0,
\]
so $\theta_1>\pi/24$.  The polynomial $t^3+t^2-1$ is strictly increasing for
$t>0$ and at $t=\tan(5\pi/24)$ equals
$-96-67\sqrt2+55\sqrt3+39\sqrt6>0$, so
$\theta_2<5\pi/24$.  Consequently
\begin{equation}
 2SC=\sin4\theta>\frac12
 \quad\text{throughout }(\theta_1,\theta_2).
 \label{eq:xy-middle-lower-window}
\end{equation}

\begin{theorem}[Closed XY signalling branches]
\label{thm:xy}
For $0\leq\theta\leq\theta_1$,
\begin{equation}
 \delta_A(U_\theta)=\sin4\theta,
 \qquad (q_I,q_X,q_Y,q_Z)
 =\left(C^2,\frac{S^2}{2},\frac{S^2}{2},0\right).
 \label{eq:xy-weak}
\end{equation}
For $\theta_2\leq\theta\leq\pi/4$,
\begin{equation}
 \delta_A(U_\theta)=S+\frac{S^2}{2},
 \label{eq:xy-strong}
\end{equation}
with
\begin{equation}
 (q_I,q_X,q_Y,q_Z)=\frac14
 \left(1+C^2+2C(1+S),\,S^2,\,S^2,\,
       1+C^2-2C(1+S)\right).
 \label{eq:xy-strong-channel}
\end{equation}
In particular, $\delta_A(\mathrm{iSWAP})=3/2$.
\end{theorem}

\begin{proof}
For every effective channel, the choice $A=0$ in
Eq.~\eqref{eq:xy-fixed-norm} gives
\[
 \sqrt{[w-(C^2-S^2)]^2+4S^2C^2}\geq2SC.
\]
This is the weak lower bound.  For the channel in Eq.~\eqref{eq:xy-weak}, let
$K_1=(SC/2)I-M_3(A)$.  Its leading principal minor is
$S(C-2AS)/2$, and its $(2,3)$ principal minor is
$AC^2(A-1)(C-1)(C+1)$.  The other two order-two principal
minors coincide and equal
\[
 \left(\frac{SC}{2}-AS^2\right)\frac{SC}{2}
 -\frac{A(1/2-A)}{2}\left[C^2(1-C)^2+S^2\right].
\]
As a quadratic in $A$, its leading coefficient is
$(C-1)(C^3-C^2-C-1)/2>0$.  With $\tau=\tan\theta$, its discriminant is
\[
 \frac{\tau^4}{(1+\tau^2)^8}
 \left(\tau^{12}-10\tau^{10}-8\tau^9+31\tau^8+16\tau^7
 -36\tau^6-32\tau^5+23\tau^4+16\tau^3
 +14\tau^2+8\tau-7\right).
\]
The first threshold obeys $\tan\theta_1<27/100$; this rational enclosure is
established in Appendix~\ref{app:xy-audit} from the strict sign change
$p_1(1/4)<0<p_1(27/100)$ of the exact threshold polynomial, and it is what
makes the following estimate valid across the entire weak interval rather
than at sampled angles.  Since $0<\tau<27/100$ throughout, and since every
positive monomial in the bracket is increasing in $\tau$ on that range,
dropping the negative monomials and evaluating each positive monomial at the
\emph{upper} enclosure $\tau=27/100$ gives an upper bound for the bracket
excluding its constant term; that sum is below $3.63<7$, so the discriminant
is nonpositive.  Completing
the square therefore proves both coincident minors nonnegative.  Finally,
\[
 \det K_1=\frac{AC(C-1)}4K(A),
\]
where $K$ is convex and
\[
 K(0)=S[2SC-(3C^3+C^2-C-1)],\qquad
 K(1/2)=-C(C+1)S(C-S).
\]
Both endpoint values are nonpositive precisely through the weak interval;
convexity gives $K(A)\leq0$.  Moreover,
\[
 \det M_3=-\frac{AC^2(2A-1)^2(C-1)^3(C+1)}4\geq0,
\]
while the sum of its order-two principal minors is nonpositive.  Thus $M_3$
has one nonnegative and two nonpositive eigenvalues: three positive
eigenvalues would make the order-two sum positive, while two positive and one
negative would make the determinant negative; the zero cases follow by
continuity.  The principal minors
above give $K_1\geq0$, hence $\lambda_{\max}(M_3)\leq SC/2$; together with the
inertia identity following Eq.~\eqref{eq:xy-minmax}, this makes
Eq.~\eqref{eq:xy-fixed-norm} at most $2SC$.  This proves
Eq.~\eqref{eq:xy-weak} and identifies the first threshold.

For the strong branch, Eq.~\eqref{eq:xy-strong-channel} has
$u=C(1+S)$ and $w=C^2$.  It is CPTP exactly when
Eq.~\eqref{eq:xy-second-threshold} holds with the strong-side inequality.
At $A=1/4$, the largest eigenvalue of $M_3$ is
\[
 L=\frac{S(2+S)}8,
\]
with eigenvector proportional to $(i\sqrt2,1,1)$.  The other two eigenvalues
have sum $-S/4$ and product
$S^2(7S^2+2S-8)/64$, so they are negative throughout this regime.
Indeed $S\geq S(\theta_2)>0.96$, whereas the positive root of
$7S^2+2S-8$ is below $0.94$.  The derivatives of $L$ in both $u$ and $w$
vanish.  Since $\lambda_{\max}$ is convex for an affine Hermitian matrix
family, the fixed $A=1/4$ witness gives the global lower bound
$4L=S+S^2/2$.

It remains to prove the matching upper bound.  Put $H_3=LI-M_3(A)$.  Direct
factorization gives
\[
 \det H_3=-\frac{S^3(4A-1)^2(S+2)}{512}
 \left[A(28S^2+8S-32)+12-13S^2-4S\right].
\]
The bracket is negative at both endpoints of $A\in[0,1/2]$.  The two
nontrivial order-two principal minors are $S^2G_1/64$ and $S^2G_2/64$, where,
with $t=2A$,
\begin{align*}
 G_1={}&(16-8S^2)t^2+(6S^2-4S-16)t+(S+2)^2,\\
 G_2={}&(12S^2-16)t^2+(32-24S^2)t+13S^2+4S-12.
\end{align*}
$G_1$ has positive completed-square minimum, while $G_2$ is concave with
nonnegative endpoint values $13S^2+4S-12$ and $(S+2)^2$; the first is already
positive for $S>0.82$, below the strong-regime range.  Together with
$(H_3)_{11}=S(S+2-4AS)/8$, all principal minors are nonnegative, so
$H_3\geq0$.  In addition,
\[
 \det M_3=
 \frac{AS^3(2A-1)^2(S+2)(7S^2+2S-8)}{32}\geq0,
\]
and the sum of its order-two principal minors is nonpositive.  Thus again
$M_3$ has one nonnegative and two nonpositive eigenvalues, and
Eq.~\eqref{eq:xy-fixed-norm} is at most $4L$.  This proves
Eq.~\eqref{eq:xy-strong}.  At $S=1$, the four probabilities are $1/4$ and the
value is $3/2$.
\end{proof}

Before treating the intermediate regime we isolate its convex-analytic
skeleton, so that the algebraic elimination that follows can be carried out
against a single, already-justified optimality certificate.

\begin{lemma}[Facet selection and convex KKT sufficiency]
\label{lem:facet}
Fix $\theta$ and parametrize the axially covariant effective channels of
Eq.~\eqref{eq:xy-channel} by their Bloch eigenvalues $(u,u,w)$.  Then the
following hold.
\begin{enumerate}
\item[(i)] \emph{(Convexity.)}  The admissible set
$\mathcal T=\{(u,w): -1\leq w\leq1,\ |u|\leq(1+w)/2\}$ is a triangle with
vertices $(0,-1)$, $(1,1)$, and $(-1,1)$, and the facet $q_Z=0$ is the edge
$w=2u-1$.  The fixed-channel diamond norm
$F(u,w)=\diamondnorm{\cN-\cE_{u,w}\cT}$ is convex and continuous on
$\mathcal T$, being the norm of a map affine in $(u,w)$.
\item[(ii)] \emph{(Concavity and saddle existence.)}  For fixed $(u,w)$ the
Watrous primal value is concave in the invariant witness parameter
$A\in[0,1/2]$.  Consequently the optimization is convex--concave and its
min--max value is attained at a saddle.  On any parameter interval where both
the maximizing witness and the minimizing channel are unique, Berge's maximum
theorem makes their coordinates continuous.
\item[(iii)] \emph{(Sufficiency.)}  Suppose $(u_\star,w_\star)$ lies in the
relative interior of the edge $w=2u-1$, the directional derivative of $F$
along that edge vanishes, and $F$ admits a subgradient whose normal component
points into $q_Z\geq0$.  Then $(u_\star,w_\star)$ minimizes $F$ over the
whole triangle $\mathcal T$, not merely over the edge.
\end{enumerate}
\end{lemma}

\begin{proof}
(i) The vertex description is immediate from the constraints, and $q_Z=0$ is
exactly $w=2u-1$.  The map $\cE_{u,w}$ is affine in $(u,w)$, hence so is
$\cN-\cE_{u,w}\cT$; a norm of an affine function is convex, and convex finite
functions on a polytope are continuous.

(ii) The feasible pairs $(\rho,Y)$ of the Watrous primal form a convex set
and the objective is linear, so the value is concave in $\rho$; the invariant
density is affine in $A$, giving concavity in $A$.  A convex--concave
function on a product of compact convex sets has a saddle by Sion's minimax
theorem.  Berge's maximum theorem gives upper hemicontinuity of the optimizer
correspondences; when each is single-valued, this is ordinary continuity.

(iii) This is the standard first-order condition for a convex function on a
polytope.  Vanishing derivative along the edge together with a subgradient
whose normal component points inward means that the negative of some
subgradient lies in the outward normal cone of $\mathcal T$ at
$(u_\star,w_\star)$, i.e.\
$0\in\partial F(u_\star,w_\star)+N_{\mathcal T}(u_\star,w_\star)$.  For a
convex $F$ this normal-cone inclusion is sufficient, not merely necessary,
for a global minimum over $\mathcal T$.  In particular no other edge or
vertex of the triangle can beat it, so the remaining facets need not be
examined separately.
\end{proof}

Lemma~\ref{lem:facet}(iii) is what reduces the intermediate regime to a
single algebraic computation: it suffices to exhibit one facet point with
vanishing tangential derivative and a positive normal multiplier.

\begin{theorem}[Intermediate XY branch]
\label{thm:xy-middle}
For $\theta_1<\theta<\theta_2$, put $t=\tan\theta$.  Then
$\delta_A(U_\theta)$ is the unique root $\Delta\in(1/2,2)$ of
\begin{align}
 Q_+(\Delta,t)={}&
 (t^6+t^5+3t^4+2t^3+3t^2+t+1)\Delta^4\nonumber\\
 &+(2t^7-46t^6-122t^5-170t^4-138t^3-58t^2-14t+2)\Delta^3\nonumber\\
 &+(-84t^7+24t^6+72t^5+96t^4-68t^3-56t^2-32t)\Delta^2\nonumber\\
 &+(96t^7+64t^6+160t^5+384t^4+320t^3+128t^2)\Delta\nonumber\\
 &+64t^5(2t^2+2t+1).
 \label{eq:xy-middle-quartic}
\end{align}
\end{theorem}

\begin{remark}[Computer-assisted status]
\label{rem:cas}
Theorem~\ref{thm:xy-middle} is a computer-assisted result.  The convex
analysis selecting the active facet, recorded in Lemma~\ref{lem:facet} and in
the first part of the proof below, is fully human-checkable.  The subsequent
eliminations are not: the exclusion of the conjugate sheet $W_-$, the
resultant factorization producing $Q_+$, the positivity of its cofactors, and
the Sturm root count are polynomial identities and sign conditions in up to
four variables with integer coefficients of large height.  They are stated in
Appendix~\ref{app:xy-audit} and \emph{verified in exact rational arithmetic}
by the script \texttt{tools/derive\_xy\_middle.py} accompanying this
manuscript, which is part of the proof of record rather than an illustration.
Because every such step is an identity or a strict inequality between
polynomials with exact rational coefficients, the verification is decidable
and involves no floating-point arithmetic, rounding, or numerical root
finding.  A reader wishing to check Theorem~\ref{thm:xy-middle} independently
should re-run that script, or re-derive the same resultants and Sturm
sequences in any computer algebra system.
\end{remark}

\begin{proof}
We first prove the active facet.  On $q_Z=0$, write $w=2u-1$,
$B=1/2-A$, and let $\ell$ be the largest eigenvalue of $M_3(A)$.  Put
\[
 m=(u-C)^2+S^2,\quad N=(u-C^2)^2+C^2S^2,\quad
 \Theta=u(1-u)(1-C)^2.
\]
The characteristic polynomial becomes
\begin{equation}
 P(\ell;A,u)=\ell^3-A(1-u)\ell^2
 -(ABm+B^2N)\ell-AB^2\Theta.
 \label{eq:xy-middle-characteristic}
\end{equation}
Thus $\det M_3=AB^2\Theta\geq0$, while the sum of its order-two
principal minors is $-(ABm+B^2N)\leq0$.  In the interior $M_3$ has one
positive and two negative eigenvalues, so the fixed-witness objective is
$4\ell$ and $P_\ell>0$.

At $A=0$, the largest root is $\ell_0=\sqrt N/2$.  Implicit differentiation
of Eq.~\eqref{eq:xy-middle-characteristic}, using
$P_\ell(\ell_0;0,u)=N/2$, gives
\begin{equation}
 4\ell'(0)=\frac{2}{N}\Xi(u),\qquad
 \Xi(u)=(1-u)N+\Theta-\sqrt N(2N-m).
 \label{eq:xy-activity}
\end{equation}
Thus $\Xi$ is the exact one-sided activity test for moving from $A=0$ into
$A>0$.

For fixed channel parameters, the diamond-SDP value is concave in its input
density, and the resulting facet objective $f(u)$ is convex and continuous,
by Lemma~\ref{lem:facet}(i)--(ii).  In the interior the maximizing $A$ is
unique: concavity makes its maximizer set an interval, while a nontrivial
interval would make the cubic $P(\ell;A,u)$ vanish identically in $A$ despite
its nonzero leading coefficient $-\Theta$.
If its unique maximizing witness were $A=0$, Danskin's theorem would give
\[
 f'(u)=\frac{2(u-C^2)}
 {\sqrt{(u-C^2)^2+C^2S^2}},
\]
so an interior channel minimizer could occur only at $u=C^2$.  At that point
Eq.~\eqref{eq:xy-activity} reduces to
\[
 \Xi(C^2)=
 C(1-C)S[2SC-(3C^3+C^2-C-1)],
\]
which is positive after $\theta_1$; hence $A=0$ is not maximizing there.
The endpoint derivatives in $u$ point inward, while the derivative toward
smaller $A$ excludes $A=1/2$.  Thus the middle saddle has
$0<u<1$ and $0<A<1/2$.

A flat interval of minimizing $u$ values is impossible.  Indeed,
\[
 \operatorname{disc}_A P
 =\frac{(C-1)^2\ell^2(2\ell+u)}{16}\mathcal R(u,\ell,C),
 \qquad [u^7]\mathcal R=(1-C)^2-8\ell.
\]
If $A=0$ on a subinterval, the boundary value
$2\sqrt{(u-C^2)^2+C^2S^2}$ is strictly convex and cannot be constant.
Hence any flat interval contains an open subinterval with $0<A<1/2$, where
the displayed discriminant applies.
At every point of that subinterval, witness stationarity gives
$P=P_A=0$, so eliminating $A$ makes
$\operatorname{disc}_A P$ vanish at the same fixed value of $\ell$.
Constancy on an interval would therefore force
$\mathcal R(u,\ell,C)$ to vanish identically as a polynomial in $u$, hence
$4\ell=(1-C)^2/2<1/2$, whereas the $A=0$ witness gives
$4\ell\geq2SC>1/2$ by
Eq.~\eqref{eq:xy-middle-lower-window}.  Both optimizer coordinates are
therefore unique, and Lemma~\ref{lem:facet}(ii) makes the facet saddle
continuous in $\theta$.

It remains to check the normal KKT multiplier.  Holding $w$ fixed before
restricting to the facet gives
\[
 P_u=2ABG,\qquad
 G=(C-u)\ell+B\Psi,\qquad
 \Psi=u^2-C(1+C)u+C.
\]
Here $\Psi>0$ because its discriminant is
$C(C-1)(C^2+3C+4)<0$.  Eliminating $A,u,\ell$ from
$P=P_A=P_u+2P_w=G=0$, after saturating the interior factors, gives
\[
 5C^3+C^2+3C-1=0.
\]
This is exactly the second-threshold equation.  At $\theta_1$ the
limiting weak saddle has $A=0$, $u=C^2$, $\ell=SC/2$, and
\[
 G=\frac{C(1-C)}2[(1+C)+SC]>0.
\]
Therefore $G>0$ throughout $(\theta_1,\theta_2)$: by
the continuity just proved, the facet saddle varies continuously in $\theta$,
so $G$ can change sign on the open interval only by vanishing there, and the
elimination above shows that this happens only at $\theta_2$.  The gradient
of the convex channel objective is then a positive multiple of the inward
normal to $q_Z\geq0$: tangential stationarity gives zero directional
derivative along the facet, while $G>0$ supplies a subgradient with inward
normal component, whose negative lies in the outward normal cone.
Lemma~\ref{lem:facet}(iii) applies verbatim, so the facet point is the global
minimizer over the entire CPTP triangle, and no other edge or vertex needs to
be examined.
Consequently $\Delta=4\ell=\delta_A(U_\theta)$.  The $A=0$ witness and
Eq.~\eqref{eq:xy-middle-lower-window} give
$\Delta\geq2SC>1/2$, while
Eq.~\eqref{eq:universal-swap-bound} gives $\Delta\leq3/2<2$.  This is the
root window used below.

It remains to eliminate the unique saddle and identify the physical algebraic
sheet.  After putting $a=2A$ and $\Delta=4\ell$, the stationary characteristic
polynomial is quadratic in $u$.  Its discriminant splits into two conjugate
factors $W_+(a,\Delta,\pm t)$.  An exact rational lower bound excludes the
minus sheet throughout an enclosing rational window for the middle interval.
Eliminating $a$ from the remaining sheet and witness stationarity gives
$Q_+(\Delta,t)$ times factors that are strictly positive on the physical
domain.  Finally, the discriminant of $Q_+$ is nonzero on that window, and an
exact Sturm count gives one root in $(1/2,2)$.  Substitution of the two closed
branches proves exact joining at both thresholds.  All factors, interval
bounds, and saturation steps are given in
Appendix~\ref{app:xy-audit}; no floating-point root selection enters this
argument.  As recorded in Remark~\ref{rem:cas}, these four eliminations are
carried out and checked in exact rational arithmetic by the accompanying
script rather than by hand.

Finally, we record in what sense this value is certified from below.  In the
weak and strong regimes the lower bound is an explicitly exhibited input
state.  Here, by contrast, the certificate is the convex--concave saddle of
Lemma~\ref{lem:facet}(ii): the value is attained at the interior invariant
witness $\rho(A_*)$ of the unique saddle, and Sion's minimax theorem
guarantees that no effective channel can do better against it.  The witness is
therefore operational and matching, but it is specified algebraically as the
saddle coordinate rather than in closed form.
\end{proof}

\begin{figure}[t]
\centering
\begin{tikzpicture}
\begin{axis}[
 width=0.82\textwidth,
 height=0.42\textwidth,
 xlabel={XY interaction angle $\theta$},
 ylabel={optimizer parameter},
 xmin=0, xmax=0.7853982,
 ymin=-1.02, ymax=1.05,
 xtick={0,0.25827052,0.64661551,0.78539816},
 xticklabels={$0$,$\theta_1$,$\theta_2$,$\pi/4$},
 legend style={at={(0.02,0.04)},anchor=south west},
 grid=major,
]
\addplot[blue!65!black,very thick]
 table[x=theta,y=transverse] {xy_phase.dat};
\addlegendentry{transverse eigenvalue $u_*$}
\addplot[green!45!black,very thick,dashed]
 table[x=theta,y=longitudinal] {xy_phase.dat};
\addlegendentry{longitudinal eigenvalue $w_*$}
\addplot[orange!85!black,very thick,dashdotted]
 table[x=theta,y=even_weight] {xy_phase.dat};
\addlegendentry{witness weight $A_*$}
\draw[densely dashed] (axis cs:0.25827052,-1.02)
 -- (axis cs:0.25827052,1.05);
\draw[densely dashed] (axis cs:0.64661551,-1.02)
 -- (axis cs:0.64661551,1.05);
\end{axis}
\end{tikzpicture}
\caption{Optimal effective-channel eigenvalues and invariant witness along the
XY line.  The weak regime has $A_*=0$; the middle saddle moves along the CPTP
facet $w_*=2u_*-1$ with $0<A_*<1/2$; and the strong regime locks to
$A_*=1/4$.  The middle values are a numerical reconstruction of the unique
algebraic saddle proved in Theorem~\ref{thm:xy-middle}, included for
visualization rather than as proof.}
\label{fig:xy-phase}
\end{figure}

The full curve satisfies
$\delta_A(\theta)=\delta_A(\pi/2-\theta)$ and returns to zero at
$\theta=\pi/2$, where $U_{\pi/2}=Z\otimes Z$.  The symmetry follows from
$U_{\pi/2-\theta}=(Z\otimes Z)\overline{U_\theta}$, together with invariance
of the diamond norm under unitary conjugation and transposition.

\section{Discussion}

The Ising family provides the baseline mechanism: one dephasing channel is
optimal for every coupling, and two product inputs already witness the exact
defect.  The main transitions occur in the other two families.  Along partial
SWAP, the optimal effective channel itself changes: it remains the identity until
$\sin\phi=(d^2-3)/(d^2-1)$, then moves into the depolarizing interval and
reaches complete depolarization at SWAP.  The optimal diamond witness likewise
generally uses a reference system, as encoded by a purification of the
symmetric/antisymmetric density in
Lemma~\ref{lem:fixed}; Fig.~\ref{fig:partial-swap-phase} displays both
optimizer trajectories for qubits.

The XY line has less symmetry and a genuine active-set phase diagram.  Its weak
branch is controlled by the one-excitation witness, its central iSWAP branch by
the maximally mixed witness, and the intervening interval moves through an
algebraic saddle, as shown in Fig.~\ref{fig:xy-phase}.  The iSWAP value $3/2$
equals the qubit SWAP value despite the
gates belonging to different local-equivalence classes in the canonical
classification of two-qubit gates~\cite{zhang2003}.  By the universal
bound~\eqref{eq:universal-swap-bound}, which applies to every unitary on
$\mathbb C^2\otimes\mathbb C^2$, iSWAP is therefore a maximally
signalling two-qubit unitary, not merely a gate that happens to share SWAP's
value.

The threshold is operational rather than spectral: it is where the outer
channel optimization first benefits from adding depolarization.  It is not
visible from operator entanglement alone.  This supplies a concrete example in
which a coarse-grained autonomy question differs from a Hilbert--Schmidt
interaction diagnostic.  For repeated-interaction and homogenization models,
the threshold also separates a regime in which the incoming visible state is
the best autonomous predictor from one in which deliberately forgetting part
of that state improves the worst-case prediction.  The defect is a worst-case
channel-simulation error, however, not a communication capacity or an
achievable signalling rate.

For example, in a qubit collision with $\sin\phi=4/5$, retaining the incoming
visible state corresponds to $\lambda=1$ and has worst-case error
$2\sin\phi=1.6$.  The physical root of
Eqs.~\eqref{eq:qubit-swap-stationarity}--\eqref{eq:qubit-swap-quartic} is
$\lambda_*=0.339610\ldots$, for which the error is
$1.415252\ldots$.  Thus the minimax one-step predictor keeps only about one
third of the incoming Bloch vector; this is the operational gain represented
by the noise-activation transition.

The result is limited in three important ways.  First, Eq.~\eqref{eq:defect}
is a pre-existing signalling measure, not a newly defined quantity.  Second,
the exact CNOT and SWAP endpoints are prior results
\cite{barsse2025}; the contribution here is the continuous-family evaluation.
Third, for a generic two-qubit interaction in the canonical Cartan
parametrization~\cite{zhang2003},
$e^{-i(\alpha XX+\beta YY+\gamma ZZ)}$, a closed formula remains open.
Joint-Pauli symmetry
reduces the latter problem to optimization over the Pauli-channel tetrahedron, but
the three leakage sectors interfere in trace norm and the partial-SWAP
$U(2)$ symmetry is absent.
In addition, the strong partial-SWAP regime is solved as an explicit scalar
convex minimization rather than by a generic radicals formula for its
minimizer.

\section{Reproducibility}

The accompanying repository separates three levels of checking, which differ
in evidential status and should not be conflated.  Level one is a
dependency-free Gaussian-rational suite that verifies residual identities,
physical witnesses, channel constraints, boundary values, and exact
substitutions into the stated formulas.  These are regression certificates:
they do not independently prove the general block spectra, and no theorem
depends on them.
The focused suite is
run from the repository root with
\begin{verbatim}
python3 -m unittest discover -s tests -p 'test_aqc_*.py' -v
\end{verbatim}
The larger symbolic identities in Theorem~\ref{thm:xy-middle} are level two.
As stated in Remark~\ref{rem:cas}, these are \emph{not} optional: the sheet
exclusion, resultant factorization, cofactor positivity, and Sturm count that
complete the proof of Theorem~\ref{thm:xy-middle} are performed and checked in
exact rational arithmetic by the first script below, which is part of the
proof of record.  The word ``optional'' applies only to the extra Python
dependency, not to the evidential role.  The second script is an independent
reconstruction, in the sense of Remark~\ref{rem:generald}.
\begin{verbatim}
python3 -m pip install -e '.[proof]'
python3 tools/derive_xy_middle.py
python3 tools/verify_block_reductions.py
\end{verbatim}
The second script constructs the unreduced physical residual Choi operators
directly and checks their characteristic polynomials against the
partial-SWAP spectrum for $d=2,3,4$ and the XY scalar-plus-$M_3$ reduction.
It does not import the certificate package.
The defect and optimizer phase-diagram data are generated without third-party
Python dependencies:
\begin{verbatim}
python3 paper/generate_xy_plot.py
\end{verbatim}
The code is a certificate and regression harness, with the single exception
recorded in Remark~\ref{rem:cas}: the analytic proofs do not depend on
floating-point optimization anywhere, and they depend on computer algebra only
for the exact polynomial eliminations behind Theorem~\ref{thm:xy-middle}.
The physical-Choi reconstruction
checks representative dimensions and parameters, while the all-$d$ statement
is the analytic block calculation in Lemma~\ref{lem:fixed}.  The generic
interior strong-branch minimizer is characterized by
Eqs.~\eqref{eq:scalar-min}--\eqref{eq:qubit-swap-quartic}, rather than encoded
as a rational test value.

As an independent implementation check, the repository also contains
\begin{verbatim}
python3 -m pip install -e '.[numerical]'
python3 tools/check_numerical_sdp.py
\end{verbatim}
This optional CVXPY program does not import the exact certificate package.  It
optimizes jointly over an unrestricted CPTP effective channel and the dual
diamond-norm SDP, then compares the unreduced numerical optimum with sample
points in every analytic regime.  These floating-point checks are not used as
proof.

\subsection*{Data and code availability}

This work generates no experimental data.  All figure data files in the
\texttt{paper/} directory are produced deterministically from
\texttt{generate\_xy\_plot.py} using only the Python standard library, and are
regenerated and compared byte-for-byte in continuous integration.  The
manuscript source, the exact-arithmetic certificate package, the symbolic
proof scripts of Remark~\ref{rem:cas}, and the independent numerical
cross-check are openly available with pinned dependency versions at
\begin{center}
\url{https://github.com/richorama/quantum-directional-signalling}.
\end{center}
The archived release DOI and immutable version identifier should be taken from
the repository's \texttt{CITATION.cff} file for the submitted version.
Software is released
under the MIT licence; the manuscript, figures, and data files are released
under CC~BY~4.0.

\subsection*{Use of AI tools}

The work was developed with the assistance of AI tools for algebraic
exploration, drafting, and code generation.  Every theorem statement, proof,
and certificate was checked by the author, who is solely responsible for all
claims and errors.

\section{Conclusion}

Directional signalling can be evaluated sharply beyond isolated gates.  The
Ising formula $|\sin2\theta|$ supplies a transparent baseline and continuous
extension of the known CNOT point.  The principal results are the two
nontrivial optimizer transitions.  In every finite dimension, partial SWAP
reduces to an explicit scalar convex minimization with
weak branch $2\sin\phi$, transition
$(d^2-3)/(d^2-1)$, and known SWAP endpoint $2(1-1/d^2)$.  The calculation
of the XY line gives an exact three-regime curve, including the apparently new
result that iSWAP reaches the global two-qubit signalling ceiling $3/2$, and a
unique quartic root in the intermediate regime.  These
results provide controlled
settings in which an optimal
autonomous effective channel changes nonperturbatively with interaction
strength.  The axially symmetric XXZ surface is the next tractable step toward
the full Cartan chamber~\cite{zhang2003}.

\appendix
\section{Prior-art search and attribution boundary}
\label{app:literature}

The novelty assessment was repeated on 30 July 2026 using the primary text and
metadata of Refs.~\cite{barsse2024,barsse2025}, arXiv title/abstract searches,
author searches for the four authors of those papers, and forward-citation
queries.  Search terms were combined across four groups:
\begin{enumerate}
 \item directional signalling, causal influence, semicausal and approximately
 semicausal channels;
 \item diamond distance to reduced, autonomous, or simulated subsystem
 dynamics;
 \item partial SWAP, collision models, homogenization, Ising, XY, and iSWAP;
 \item Cartan two-qubit gates, operator entanglement, tensor-product-structure
 distance, and quantum model reduction.
\end{enumerate}
No independent paper was found that evaluates
Eq.~\eqref{eq:defect} continuously on the Ising, partial-SWAP, or XY/iSWAP
families.  Searches for close phrases such as ``partial SWAP'' with ``diamond
norm'', ``iSWAP'' with ``diamond norm'', and ``semicausal'' with ``diamond
norm'' found no overlapping calculation.  Citation searches for
Refs.~\cite{barsse2024,barsse2025} likewise found no follow-up evaluating these
families.

This negative search result is not a proof of novelty.  Title and abstract
indexes can miss results stated only in the body of a paper, and terminology
such as autonomy, intertwining defect, model reduction, or approximate
semicausality is not standardized.  The precise attribution boundary used
here is therefore conservative:
\begin{itemize}
 \item the signalling quantity and its relation to causal influence are from
 Ref.~\cite{barsse2024};
 \item the exact CNOT value, exact arbitrary-dimensional SWAP value,
 depolarizing SWAP optimizer, and equal-dimension universal upper bound are
 from Ref.~\cite{barsse2025};
 \item the continuous-family formulas, thresholds, effective channels, and
 matching witnesses proved here are the claimed contributions.
\end{itemize}
The repository records the search scope and reproducible claim audit so that
this assessment can be updated if closer prior art is identified.

\section{Audit details for the intermediate XY branch}
\label{app:xy-audit}

This appendix records the boundary and saturation details used in
Theorem~\ref{thm:xy-middle}.  They are included to make the exact eliminations
auditable without reconstructing the branch selection from numerical output.

\subsection{Boundary exclusion and uniqueness}

On the facet $w=2u-1$, Eq.~\eqref{eq:xy-middle-characteristic} satisfies
\[
 P_A\big|_{A=1/2,\,\ell=(1-u)/2}
 =\frac{u(1-u)(1-C)}2>0.
\]
Since $P_\ell>0$, the largest eigenvalue decreases as $A$ approaches
$1/2$; that endpoint cannot maximize the witness objective.  At $A=0$, the
objective is
\[
 2\sqrt{(u-C^2)^2+C^2S^2}.
\]
The maximizing $A$ is unique by the argument in the main proof.  Therefore,
if $A=0$ is active, Danskin's theorem identifies the derivative of the
maximized facet objective with the derivative of this displayed boundary
function.  It vanishes only at $u=C^2$.  The derivative into $A>0$ there has
the sign of $\Xi(C^2)$ from Eq.~\eqref{eq:xy-activity}, namely
\[
 C(1-C)S[2SC-(3C^3+C^2-C-1)],
\]
which is positive exactly after the first threshold, so $A=0$ is not active
at the only possible stationary channel value.  Thus the middle saddle has
$0<A<1/2$.  Direct endpoint evaluation gives
$f'(0^+)=-2$ and $f'(1^-)=2S$, so also $0<u<1$.

For fixed $u$, concavity in $A$ makes the maximizer set an interval.  Any
nontrivial interval would make the cubic
$P(\ell;A,u)$ vanish identically in $A$, contradicting the coefficient
$-\Theta$ of $A^3$.  For fixed $C$, convexity in $u$ makes the minimizer set
an interval.  A flat portion on the $A=0$ boundary is excluded by strict
convexity of the displayed square root.  In the interior it would force the
$A$-discriminant polynomial $\mathcal R$ in the main proof to vanish identically in
$u$, including its leading coefficient $(1-C)^2-8\ell$.  This contradicts
$4\ell\geq2SC>1/2$ from
Eq.~\eqref{eq:xy-middle-lower-window}.  Hence both optimizers are unique, and Berge's theorem
gives their continuity in $C$.

\subsection{Multiplier elimination and saturation ledger}

The zero-multiplier calculation starts from the unrestricted characteristic
polynomial and the equations
\[
 P=0,\qquad P_A=0,\qquad P_u+2P_w=0,\qquad G=0.
\]
The physical interior permits saturation by
\[
 A,\ B,\ u,\ 1-u,\ C,\ 1-C^2,\ \ell,\ P_\ell,\ \Psi.
\]
Every factor is nonzero: the boundary factors were excluded above;
$0<C<1$ and $\ell>0$ hold in the middle interval;
$P_\ell>0$ follows from the simple positive eigenvalue; and
$\Psi>0$ follows from its negative discriminant.  The substitution
$B=(u-C)\ell/\Psi$ therefore loses no physical solution.  The remaining
generic component $u=-2\ell$ is incompatible with $u,\ell>0$.
After these exclusions, the saturated elimination ideal in
$\mathbb{Q}[C]$ is generated by $5C^3+C^2+3C-1$.  It vanishes only at
$\theta_2$, establishing the nonvanishing multiplier used in the KKT
argument.  The optional symbolic script also prints the unsaturated resultant
factors and the justification for excluding each one.

\subsection{Rational enclosure of the threshold interval}

The elimination bounds below use a rational window rather than decimal
approximations.  In the tangent variable, the first and second threshold
polynomials are
\[
 p_1(t)=t^6-2t^5-3t^4+7t^2+2t-1,\qquad
 p_2(t)=t^3+t^2-1.
\]
Exact evaluation gives
\[
 p_1(1/4)=-\frac{311}{4096}<0,\qquad
 p_1(27/100)=\frac{31874409089}{10^{12}}>0,
\]
and
\[
 p_2(3/4)=-\frac1{64}<0,\qquad
 p_2(19/25)=\frac{259}{15625}>0.
\]
Using the uniqueness of the positive threshold roots therefore gives
\[
 \frac14<\tan\theta_1<\frac{27}{100}
 <\frac34<\tan\theta_2<\frac{19}{25}.
\]
This both justifies the enclosing interval $[1/4,19/25]$ and supplies the
$t<27/100$ bound used in the weak-branch principal-minor estimate.

\subsection{Physical sheet, resultant, and root isolation}

Put $a=2A$ and
$F=64(1+t^2)^2P(\Delta/4;a/2,u)=P_2u^2+P_1u+P_0$.
At the saddle $F_u=0$, and direct factorization gives
\[
 \operatorname{disc}_uF=4W_+(a,\Delta,t)W_+(a,\Delta,-t),
\]
where
\[
 W_+=(1+t^2)^2\left[
 a\left(\Delta+\frac{4(1-a)t^2}{1+t^2}\right)^2-2\Delta^2
 +\frac{8\Delta(1-a)t}{1+t^2}
 \left(\frac{2at^2+1-t^2}{1+t^2}\right)\right].
\]
With $\sigma=t^2/(1+t^2)$ and $b=1-a$, the conjugate sheet obeys the exact
rearrangement
\[
 -\frac{W_+(a,\Delta,-t)}{(1+t^2)^2}
 =(1+b)\Delta^2+4b\Delta\kappa-16(1-b)b^2\sigma^2,
\]
where
\[
 \kappa=SC-2a\sigma(1-S)\geq S+C-1>0.
\]
Using $a(1-a)\leq1/4$, $\Delta\geq1/2$, and
$t\in[1/4,19/25]$ gives the rational lower bound
$171381/972196$.  Thus the conjugate sheet cannot contain the physical
saddle.

The second stationary equation, after using $F_u=0$, is
$G_8(a,\Delta,t)=0$, defined without expansion by
\[
 -16G_8=(\partial_aP_2)P_1^2
 -2P_1(\partial_aP_1)P_2+4(\partial_aP_0)P_2^2.
\]
Exact elimination gives
\begin{align*}
 \operatorname{Res}_a(W_+,G_8)
 ={}&-2^{36}\Delta^{12}t^{28}(1+t^2)^{10}E_1E_2^2Q_+(\Delta,t),\\
 E_1={}&\Delta^2(1+3t^2)+4(1-t^2)(1+\Delta),\\
 E_2={}&\Delta(1+t^2)^2+2(1+t)^2(1+2t-t^2).
\end{align*}
Both $E_1$ and $E_2$ are strictly positive for $0<t<1$ and
$\Delta>0$, so the physical saddle satisfies $Q_+(\Delta,t)=0$.

For root isolation,
\begin{align*}
 Q_+(1/2,t)&=\frac1{16}
 (2484t^7+2565t^6+2349t^5+3119t^4+2014t^3
 +687t^2-155t+5)>0,\\
 Q_+(2,t)&=-32(9t^5+5t^4+22t^3+12t^2+7t-1)<0
\end{align*}
on $[1/4,19/25]$, while
\[
 \operatorname{disc}_\Delta Q_+
 =2^{16}t^7(1+t^2)(t^3-t^2-t-1)^2N_9(t)^3
\]
with
\[
 N_9(t)=27t^9+378t^8+1026t^7+2106t^6+2772t^5
 +2646t^4+1658t^3+678t^2+141t+16.
\]
The leading coefficient $(1+t^2)^2(t^2+t+1)$ and the discriminant are
nonzero there, preventing roots from entering, leaving, colliding, or
escaping to infinity.  The exact Sturm variation difference at $t=1/2$ is
one between $\Delta=1/2$ and $\Delta=2$.

Finally, substituting the weak and strong values into $Q_+$ produces,
respectively, the squared threshold factors
\[
 p_1(t)^2,\qquad p_2(t)^2,
\]
with strictly positive cofactors on the enclosing window.  Hence the unique
physical root joins both closed branches exactly.

\begin{thebibliography}{99}

\bibitem{barsse2024}
K.~Barsse, P.~Perinotti, A.~Tosini, and L.~Vaglini,
``Causal influence versus signalling for interacting quantum channels,''
\emph{Phys. Rev. Research} \textbf{6}, 043305 (2024),
\href{https://arxiv.org/abs/2309.07771}{arXiv:2309.07771}.

\bibitem{barsse2025}
K.~Barsse, P.~Perinotti, A.~Tosini, and L.~Vaglini,
``Disentangling signalling and causal influence,'' preprint (2025),
\href{https://arxiv.org/abs/2505.14120}{arXiv:2505.14120}.

\bibitem{beckman2001}
D.~Beckman, D.~Gottesman, M.~A. Nielsen, and J.~Preskill,
``Causal and localizable quantum operations,''
\emph{Phys. Rev. A} \textbf{64}, 052309 (2001),
\href{https://arxiv.org/abs/quant-ph/0102043}{arXiv:quant-ph/0102043}.

\bibitem{eggeling2002}
T.~Eggeling, D.~Schlingemann, and R.~F. Werner,
``Semicausal operations are semilocalizable,''
\emph{Europhys. Lett.} \textbf{57}, 782 (2002),
\href{https://arxiv.org/abs/quant-ph/0104027}{arXiv:quant-ph/0104027}.

\bibitem{ksw2008}
D.~Kretschmann, D.~Schlingemann, and R.~F. Werner,
``A continuity theorem for Stinespring's dilation,''
\emph{J. Funct. Anal.} \textbf{255}, 1889 (2008),
\href{https://arxiv.org/abs/0710.2495}{arXiv:0710.2495}.

\bibitem{watrous2009}
J.~Watrous, ``Semidefinite programs for completely bounded norms,''
\emph{Theory Comput.} \textbf{5}, 217 (2009).

\bibitem{watrous2018}
J.~Watrous, \emph{The Theory of Quantum Information}
(Cambridge University Press, Cambridge, 2018).

\bibitem{zanardi2001}
P.~Zanardi, ``Entanglement of quantum evolutions,''
\emph{Phys. Rev. A} \textbf{63}, 040304(R) (2001),
\href{https://arxiv.org/abs/quant-ph/0010074}{arXiv:quant-ph/0010074}.

\bibitem{beny2010}
C.~B\'eny and O.~Oreshkov,
``General conditions for approximate quantum error correction and
near-optimal recovery channels,''
\emph{Phys. Rev. Lett.} \textbf{104}, 120501 (2010),
\href{https://arxiv.org/abs/0907.5391}{arXiv:0907.5391}.

\bibitem{andreadakis2025}
F.~Andreadakis and P.~Zanardi,
``Tensor product structure geometry under unitary channels,''
\emph{Quantum} \textbf{9}, 1668 (2025),
\href{https://arxiv.org/abs/2410.02911}{arXiv:2410.02911}.

\bibitem{grigoletto2025}
T.~Grigoletto, Y.~Tao, F.~Ticozzi, and L.~Viola,
``Exact model reduction for continuous-time open quantum dynamics,''
\emph{Quantum} \textbf{9}, 1814 (2025),
\href{https://arxiv.org/abs/2412.05102}{arXiv:2412.05102}.

\bibitem{ziman2002}
M.~Ziman, P.~Stelmachovic, V.~Buzek, M.~Hillery, V.~Scarani, and N.~Gisin,
``Quantum homogenization,''
\emph{Phys. Rev. A} \textbf{65}, 042105 (2002),
\href{https://arxiv.org/abs/quant-ph/0110164}{arXiv:quant-ph/0110164}.

\bibitem{scarani2002}
V.~Scarani, M.~Ziman, P.~Stelmachovic, N.~Gisin, and V.~Buzek,
``Thermalizing quantum machines: Dissipation and entanglement,''
\emph{Phys. Rev. Lett.} \textbf{88}, 097905 (2002),
\href{https://arxiv.org/abs/quant-ph/0110088}{arXiv:quant-ph/0110088}.

\bibitem{ciccarello2022}
F.~Ciccarello, S.~Lorenzo, V.~Giovannetti, and G.~M. Palma,
``Quantum collision models: Open system dynamics from repeated
interactions,''
\emph{Phys. Rep.} \textbf{954}, 1 (2022),
\href{https://arxiv.org/abs/2106.11974}{arXiv:2106.11974}.

\bibitem{zhang2003}
J.~Zhang, J.~Vala, S.~Sastry, and K.~B. Whaley,
``Geometric theory of nonlocal two-qubit operations,''
\emph{Phys. Rev. A} \textbf{67}, 042313 (2003),
\href{https://arxiv.org/abs/quant-ph/0209120}{arXiv:quant-ph/0209120}.

\bibitem{zanardi2024}
P.~Zanardi, N.~Dallas, F.~Andreadakis, and S.~Lloyd,
``Operational quantum mereology and minimal scrambling,''
\emph{Quantum} \textbf{8}, 1406 (2024),
\href{https://arxiv.org/abs/2212.14340}{arXiv:2212.14340}.

\end{thebibliography}

\end{document}
